\chapter{Path to an Integrated Circuit}
\label{ch:ch6_asic}

% [DRAFT v1] Source material: docs/sintesis_resultados.md, docs/via_asic.md,
% synth/sintetiza.py, results/sintesis_sg13g2.csv.

The engine is designed to end up as a dedicated integrated circuit that
serves keys to a microcontroller, and the FPGA is its verification
vehicle, not its destination. This chapter takes the same VHDL that
passes the testbenches of Chapter~\ref{ch:ch5_verification}, synthesizes
it on an open standard-cell library, and states what a fabrication run
would cost and what would have to change in the entropy source. A
tape-out is explicitly out of scope of this thesis --- not for cost or
area, as will be seen, but for schedule --- and everything here is
labelled as what it is: a lower bound on area from logic synthesis, with
the growth factors to core area taken from published data.

\section{Standard-cell synthesis flow}
\label{sec:synthflow}

The flow is deliberately reproducible with free tools. Yosys reads the
VHDL through the GHDL front end, synthesizes each block flattened, maps
flip-flops and logic onto the liberty file of the target library and
reports cell count and area. The library is the open \ac{PDK} of IHP,
SG13G2, a 130~nm BiCMOS process with a maintained standard-cell library,
compiled SRAM macros, I/O cells and OpenROAD support, chosen because it is
the one open \ac{PDK} that is alive, European and reachable through the
same Europractice membership the university already holds. The gate
equivalent is read from the liberty file itself rather than from memory:
$1~\mathrm{GE} = $ area of \texttt{sg13g2\_nand2\_1} $= 7.2576~\mu\mathrm{m}^2$.
No clock constraint is applied: Yosys does not use one for area, so the
figures do not depend on a target frequency. The ring oscillators are not
synthesized; they are the one block designed at transistor level
(Chapter~\ref{ch:ch3_entropy}).

\section{Area per block in gate equivalents}
\label{sec:area}

Table~\ref{tab:synth} gives the result for each block. The digital
engine --- everything except the capture buffer --- is about
\textbf{62~kGE}, and AES is 51\,\% of it: the authenticator and the
generator each instantiate their own encryption core, 15.8~kGE apiece.
The capture buffer is reported for completeness but must not be read as
part of the product: its 46.7~kGE are a 4096-bit memory mapped onto
flip-flops because the flow has no SRAM macros. In the FPGA it occupies a
block RAM at no logic cost; in silicon it would be a compiled SRAM or,
more sensibly, a much smaller buffer, because it is a test path that is
disabled at production.

\begin{table}[htbp]
\centering
\caption{Logic synthesis on IHP SG13G2, flattened, no routing.
1~GE $=7.2576~\mu\mathrm{m}^2$. Lower bound on area \simu{}.}
\label{tab:synth}
\begin{tabular}{@{}lrrrr@{}}
\toprule
Block & Cells & Flip-flops & Area ($\mu\mathrm{m}^2$) & kGE \\
\midrule
AES-128 encryption core       & 10\,439 & 262   & 114\,771 & 15.81 \\
AES-CMAC (with its AES)        & 13\,210 & 909   & 169\,766 & 23.39 \\
CTR\_DRBG with df (with its AES) & 17\,953 & 2\,355 & 267\,262 & 36.83 \\
Health tests (RCT, APT)        & 568     & 85    & 8\,986   & 1.24 \\
ERO sampler and divider        & 117     & 24    & 2\,078   & 0.29 \\
Bit clock-domain crossing      & 11      & 7     & 408      & 0.06 \\
I2C slave                      & 310     & 63    & 5\,507   & 0.76 \\
\midrule
Capture buffer (4096-bit RAM as flip-flops) & 12\,250 & 4\,140 & 339\,057 & 46.72 \\
\bottomrule
\end{tabular}
\end{table}

The numbers are internally consistent, and that consistency is how an
error in the flow was caught: before flattening, the report returned for
the authenticator and for the generator the area of their AES
sub-module alone. An authenticator weighing exactly as much as a bare AES
is impossible, so the report, not the design, was wrong. After
flattening, the authenticator adds to the AES 647 flip-flops (sub-keys,
chaining state, key) and 7.6~kGE, and the generator adds 2093 flip-flops
and 21~kGE, about half of which are its registers.

\section{Optimizations and their expected savings}
\label{sec:opt}

The synthesis is naive on purpose --- the \ac{RTL} was written to be
verifiable, not small --- and it points at three optimizations, in order
of return on effort:

\begin{enumerate}
\item \textbf{One shared AES core.} The top-level sequencer already
serializes commands: the generator and the authenticator never run at the
same time. Sharing the core saves $\approx$15.8~kGE, a quarter of the
engine, for a trivial arbiter.
\item \textbf{Compact S-box datapath.} Twenty independent S-box tables
synthesize to 15.8~kGE for the whole core; area-optimized designs in the
literature range from 2.4~kGE for an 8-bit datapath at 226 cycles per
block~\cite{moradi2011pushing} to about 5~kGE for 32-bit datapaths. Speed
is irrelevant here, so a 32-bit datapath with four shared S-boxes is the
sensible point; estimated saving, 10~kGE per core.
\item \textbf{Fewer registers in the generator.} Of its 2355 flip-flops,
the captured entropy and nonce (384), the seed material (256), the output
(256) and the temporaries of the derivation function can be reduced by
consuming the seed as a stream and sharing temporaries; roughly 800
flip-flops, about 4~kGE.
\end{enumerate}

With the three, the digital engine falls to \textbf{$\approx$20~kGE},
which is the order of magnitude the state of the art anticipated (13 to
25~kGE for engines of this kind) and which, at the published growth
factor of $\approx$2.35 from cell area to core area for an AES on this same
library, is about 0.34~mm$^2$ of core. Each of these savings is stated as
an estimate; measuring them by synthesizing the optimized variants is
listed as future work, not claimed.

\section{What the ring oscillator needs in silicon}
\label{sec:ro_silicon}

The digital blocks port with the standard flow. The entropy source does
not, and the literature that has done the crossing --- notably a 2024
study that characterizes the same ring model on an Artix-7 and on a 28~nm
FD-SOI circuit~\cite{benea_2024_flicker} --- makes three points that the
chip design must absorb.

First, the ring is rebuilt from inverter cells rather than look-up tables,
and its jitter changes in value and in character: flicker noise gains
weight in newer nodes, and the study finds, against expectation, that
counting it as a legitimate source can raise throughput for the same
entropy. The stochastic model and the estimator of
Chapter~\ref{ch:ch3_entropy} carry over unchanged; the measured $\sigma$
does not. Second, process variation makes every die different, so the
frequency and jitter of each ring must be calibrated per chip; the
published all-digital designs that survive voltage and temperature
sweeps do it with a configurable ring and a tuning loop~\cite{yang_2016_jssc}.
Third, and most important for a security device, injecting a frequency on
the supply locks the rings and removes the jitter: a 500~mV signal on
the rail reduced the key space of a commercial secure microcontroller
from $2^{64}$ to about 3300~\cite{markettos_2009_frequency}, and a few
hundred microwatts of radiated field destroyed the statistics of a
fifty-ring generator~\cite{bayon_2012_contactless}. The same calibration
loop, together with supply filtering, is the published defence, and the
health tests of Section~\ref{sec:health} are the detector of last resort.
None of this is in the FPGA prototype; it is the work that a silicon
implementation adds.

\section{Cost and schedule of an academic shuttle}
\label{sec:shuttle}

Fabrication is more accessible than the size of the word suggests. The
University of Sevilla and IMSE-CNM are members of Europractice, whose
academic price list is public: a mini@sic block in UMC~180~nm costs
3\,430~\texteuro{}, one in TSMC~65~nm 3\,691~\texteuro{}, and IHP SG13G2 is
offered at 5\,110~\texteuro{}/mm$^2$ with a free multi-project run for
open-source designs under 2~mm$^2$~\cite{europractice_prices2026}. An engine of
0.3--0.4~mm$^2$ fits many times over in the smallest block of any of
them; the chip would be limited by its pads, not its logic. The
institute's own hardware-security group has already taken ring-oscillator
primitives from Xilinx~7-series FPGAs to characterized 65~nm silicon,
which makes the continuation of this work a matter of joining an
existing line rather than opening one.

What does not fit in a master's thesis is the schedule: from tape-out to
packaged samples takes months to over a year, the price does not include
packaging, and the characterization of the returned dies is a project in
itself. For comparison, a commercial secure element such as the
ATECC608B costs 0.90~\$ in single units and ships with a NIST entropy
source validation~\cite{microchip_atecc608b_ds}; two of its direct competitors
do not mention SP~800-90 anywhere in their datasheets. The honest
position of this work is therefore not to compete on price but on
\emph{transparency of the entropy}: every bit of a key produced by this
engine can be traced to a measured jitter and a stated bound, which is
precisely what a black-box part cannot offer.
